For each fixed support, $s_n$ is bounded. The dyadic prime is unique;
each fixed odd prime has open decomposition group in the real
cyclotomic $\mathbf Z_2$-extension, since its image in
$\mathbf Z_2^\times/\{\pm1\}$ has nonzero two-adic logarithm.
Use now the finite free model and coefficient change already constructed.
Its scalar reduction is a bounded finite free $\mathbf F_2[[t]]$-complex,
and its finite cyclotomic specializations are derived quotients by
$t^{2^n}$. Smith decomposition shows that torsion disks and their Tor
terms contribute bounded dimensions, while a free cohomology generator
contributes $2^n$. The displayed dimensions therefore give generic
rank one in degree one and zero in all other degrees.

For the present curve, $E[2]$ is the direct sum of two trivial
$\mathbf F_2$-representations. Thus the residual positive complex has
rank two in degree one only. The residue field of the local ring
$B=\mathscr O[G]$ is $\mathbf F_2((t))$; the augmentation ideal of
the finite two-group lies in its maximal ideal. A minimal finite free
$B$-model has zero residual differential, so the just-computed dimensions
force exactly two generators in degree one and no others. Consequently
\[
                    C^+\otimes B\simeq B^2[-1].\tag{E.3}
\]
This uses the local ring $B$, not semisimplicity of $\mathbf F_2[G]$.

The common real sign is needed here. The real cyclotomic character
and $\varepsilon$ are trivial on complex conjugation, so the real
action is the fixed action on $T_2E\otimes\chi_0$, extended to the
group ring. Its invariant module is the fixed free rank-one
$\Z_2$-lattice extended to $\Lambda[G]$, and then to
$\mathscr O[G]$. Thus a generator can be chosen before $G$ varies.
Real $H^1$ is killed by two over these coefficient rings, so the
positive lift below costs the same single factor two for every
binary dimension.
